\chapter{Глоссарий внешних идей}
\label{ch:glossary}

Краткие описания инженерных и научных идей, на которые опирается монография, чтобы
читатель не был вынужден держать в голове внешний контекст. Порядок — тематический.

\section*{Память и активация}

\paragraph{Базовая активация ACT-R.} В когнитивной архитектуре ACT-R (Anderson)
сила следа в памяти есть функция \emph{частоты} и \emph{свежести} обращений со
степенным затуханием: $B_v=\ln\sum_k t_k^{-\alpha}$, где $t_k$ — давность $k$-го
обращения. Следствие: частые и \emph{разнесённые во времени} обращения укрепляют
след сильнее массовых (эффект распределённого повторения).

\paragraph{Spreading activation.} Модель ассоциативного припоминания
(Collins–Loftus): активация растекается от активированного узла по рёбрам, затухая
с расстоянием. Формальная реализация — случайные блуждания и персонализированный
PageRank.

\paragraph{Сети Хопфилда, обучение Хебба.} Ассоциативная память как граф со
взвешенными связями и динамикой активации, сходящейся к аттрактору (энергетический
минимум). Правило Хебба: связь между совместно активными узлами усиливается («что
вместе активируется, вместе укрепляется»). Наши fusion/консолидация — это
хеббовское обучение на графе.

\section*{Убеждения и истинность}

\paragraph{Система поддержки истинности (TMS/ATMS).} Механизм (Doyle;
de~Kleer), хранящий \emph{обоснования} выводов: для каждого производного убеждения
записано, из каких посылок оно получено. При ретракции посылки её обоснованные
следствия автоматически помечаются на пересмотр. Лечит «ползучую веру».

\paragraph{Недоопределённость (underdetermination).} Тезис Дюэма–Куайна: данные не
определяют теорию однозначно — одни и те же наблюдения совместимы с разными
гипотезами. Когнитивное проявление — преждевременный коллапс: ум схлопывает
суперпозицию совместимых гипотез в одну слишком рано.

\paragraph{Belief perseverance.} Устойчивость убеждения к опровержению его
основания: даже после опровержения посылки её следствия продолжают влиять.

\paragraph{Петлевое распространение (loopy belief propagation).} Распространение
уверенности по циклическому графу может привести к тому, что убеждение подтверждает
само себя по петле (эхо-камера), искусственно раздувая уверенность. Лечится
дисконтом по общей родословной свидетельств.

\section*{Динамика и критичность}

\paragraph{Самоорганизованная критичность (SOC).} Класс систем (Bak–Tang–Wiesenfeld;
образ — куча песка), самонастраивающихся к критической точке, где отклики следуют
степенному закону: почти все малы, изредка случаются крупные лавины. Нейронные
лавины (Beggs–Plenz) — проявление SOC в мозге.

\paragraph{Ветвящийся коэффициент $\sigma$.} Среднее число узлов, активируемых
одним активным. $\sigma<1$ — субкритично (затухает), $\sigma=1$ — критично,
$\sigma>1$ — суперкритично (runaway). Управляющий параметр лавины пропагации.

\paragraph{Гистерезис (триггер Шмитта).} Схема с двумя порогами (включения и
выключения), не дающая системе «дребезжать» у единственного порога.

\section*{Графы: спектр, поток, топология}

\paragraph{Стационарное распределение $\pi$.} Долгосрочная частота посещений узлов
случайным блужданием; глобальный объект (перронов собственный вектор матрицы
переходов), а не локальные вероятности.

\paragraph{Время перемешивания (mixing time).} Скорость сходимости блуждания к
$\pi$; управляется спектральным зазором матрицы переходов.

\paragraph{PageRank / персонализированный PageRank / RWR.} Случайное блуждание с
телепортом (рестартом): с вероятностью $\alpha$ прыжок в узел из распределения
рестарта. Телепорт делает блуждание эргодическим — гарантирует посещение всех
узлов.

\paragraph{Эффективное сопротивление $\Reff$.} Мера «электрической» удалённости
узлов, если рёбра — резисторы: $\Reff(i,j)=(e_i-e_j)^\top L^{+}(e_i-e_j)$. Большое
$\Reff$ между узлами ограничивает передачу влияния между ними.

\paragraph{Проводимость / константа Чигера.} Мера «узости» перешейков графа;
неравенство Чигера связывает её со спектральным зазором. Малая проводимость = узкое
место.

\paragraph{Over-squashing.} Удушение дальней связи: информация от экспоненциально
растущей окрестности сжимается в вектор фиксированной размерности, проходя через
узкое место графа; влияние дальних узлов затухает.

\paragraph{Over-smoothing.} Пересглаживание: при многократной диффузии
представления узлов сходятся к низкоразмерному консенсусу, теряя различимость.

\paragraph{Кривизна графа (Оливье, Форман).} Дискретные аналоги кривизны;
отрицательно искривлённые рёбра — мосты между сообществами, места over-squashing.
Полезны как локатор узких мест.

\paragraph{Анти-энтропийный скраб.} Приём из распределённых хранилищ (ZFS scrub;
Merkle-based anti-entropy gossip): фоновая случайная перепроверка всего хранилища,
поскольку неизвестно, какая запись повреждена.

\paragraph{Огрубление графа (graph coarsening), coarse-to-local.} Замена мелкого
графа крупным с сохранением релевантных свойств (спектра, разрезов, потоков).
Центральный вопрос — какой сигнал выживает при переходе от мелкого масштаба к
крупному; в памяти это консолидация.

\paragraph{Персистентный лапласиан, TDA.} Топологический анализ данных: как меняются
топологические признаки (числа Бетти, спектр лапласиана) при изменении масштаба
(фильтрации). Персистентный лапласиан — спектральное обобщение персистентной
гомологии.

\section*{Выразительность и обучение}

\paragraph{WL-тест, GIN.} Тест Вейсфейлера–Лемана — предел различающей способности
обычной передачи сообщений; GIN (Xu и др.) достигает этого предела инъективной
агрегацией. PNA — многостатистическая агрегация (сумма/среднее/макс/дисперсия).

\paragraph{Передача сообщений (GNN).} Параметрическая функция на графе, обучаемая
градиентным спуском: узел обновляется агрегированными сообщениями соседей. Наша
память — та же передача сообщений, но с локальными (не обучаемыми градиентом)
правилами.

\paragraph{LoRA.} Дешёвое дообучение: вместо всех весов обучается низкоранговая
добавка $\Delta W=BA$ (ранг $r\ll$ размерности), малая по памяти.

\paragraph{DPO (Direct Preference Optimization).} Дообучение под предпочтения без
отдельной модели вознаграждения: оптимизация на тройках (промпт, выбранный ответ,
отвергнутый ответ).

\section*{Эмоции и данные}

\paragraph{Теория оценивания (appraisal): OCC, Шерер.} Эмоция как результат оценки
события по измерениям: релевантность/новизна (ошибка предсказания), конгруэнтность
целям, валентность, интенсивность, атрибуция причины (agency).

\paragraph{Предиктивная обработка, ошибка предсказания.} Взгляд (Friston и др.), по
которому мозг минимизирует ошибку предсказания; величина ошибки — естественная мера
значимости события.

\paragraph{ConceptNet / Numberbatch.} Ассоциативно-таксономический граф здравого
смысла (взвешенные типизированные рёбра) и его родной эмбеддинг Numberbatch.

\paragraph{ATOMIC / ATOMIC10x.} Граф причинно-социального здравого смысла (если
субъект делает $Y$, то намерение/реакция/эффект). ATOMIC10x — крупная версия,
дистиллированная языковой моделью и отфильтрованная критиком.

\paragraph{Wikidata.} Фактологическая база с моделью утверждения: квалификаторы
(время), references (провенанс), ранг (недеструктивная ревизия). Ценна схемой, не
содержанием.

\paragraph{Generative Agents.} Симуляция агентов с потоком памяти, retrieval по
свежести$\times$важности$\times$релевантности и рефлексией (порождением обобщений);
прообраз субстрата.

\paragraph{Распределение Бета, сэмплирование Томпсона.} Сопряжённый апостериор для
доли успехов; удобно моделирует «успешность реакции» (успехи $\alpha$, провалы
$\beta$) и баланс исследования/эксплуатации.
